\documentclass[11pt,a4paper]{article}
\usepackage[top=20mm,bottom=20mm,left=19mm,right=19mm,headheight=15pt,headsep=6mm,footskip=10mm]{geometry}
\usepackage{fontspec}
\setmainfont{texgyrepagella}[Extension=.otf,UprightFont=*-regular,BoldFont=*-bold,ItalicFont=*-italic,BoldItalicFont=*-bolditalic]
\setsansfont{texgyreheros}[Extension=.otf,UprightFont=*-regular,BoldFont=*-bold,ItalicFont=*-italic,BoldItalicFont=*-bolditalic]
\setmonofont{lmmono10-regular.otf}[Scale=0.88]
\usepackage{amsmath,amssymb}
\usepackage{graphicx,array,booktabs,longtable,calc}
\usepackage{xcolor}
\definecolor{Ink}{HTML}{183548}
\definecolor{Accent}{HTML}{147D87}
\definecolor{Muted}{HTML}{526571}
\usepackage{microtype}
\usepackage{enumitem}
\setlist{itemsep=2pt,parsep=1pt,topsep=4pt,leftmargin=1.5em}
\usepackage{titlesec}
\titleformat{\section}{\large\sffamily\bfseries\color{Ink}}{}{0pt}{}
\titleformat{\subsection}{\normalsize\sffamily\bfseries\color{Accent}}{}{0pt}{}
\titlespacing*{\section}{0pt}{13pt}{5pt}
\titlespacing*{\subsection}{0pt}{9pt}{4pt}
\usepackage{fancyhdr,lastpage}
\pagestyle{fancy}
\fancyhf{}
\fancyhead[L]{\small\sffamily\color{Muted}SAFE FROM FINE-TUNE TO PHONE}
\fancyhead[R]{\small\sffamily\color{Muted}MAC EXPERIMENT REPORT}
\fancyfoot[L]{\footnotesize\sffamily\color{Muted}Manish | Lexsi Hackathon, Track 2}
\fancyfoot[R]{\footnotesize\sffamily\color{Muted}\thepage\ / \pageref{LastPage}}
\renewcommand{\headrulewidth}{0.3pt}
\usepackage{xurl}
\usepackage[unicode,colorlinks=true,linkcolor=Accent,urlcolor=Accent]{hyperref}
\hypersetup{pdftitle={Safe from Fine-Tune to Phone: Mac Experiment Report},pdfauthor={Manish},pdfsubject={November 4-5 execution, findings, limitations and blockers}}
\setlength{\parindent}{0pt}
\setlength{\parskip}{3.5pt}
\AtBeginDocument{\fontsize{10.5}{13}\selectfont}
\setlength{\emergencystretch}{3em}
\setlength{\LTleft}{0pt}
\setlength{\LTright}{0pt}
\setlength{\tabcolsep}{4pt}
\renewcommand{\arraystretch}{1.15}
\usepackage{etoolbox}
\AtBeginEnvironment{longtable}{\fontsize{9.3}{11.7}\selectfont}
\AtBeginEnvironment{enumerate}{\interlinepenalty=10000}
\setcounter{secnumdepth}{0}
\widowpenalty=10000
\clubpenalty=10000
\raggedbottom
\providecommand{\tightlist}{\setlength{\itemsep}{1pt}\setlength{\parskip}{0pt}}
\providecommand{\pandocbounded}[1]{#1}
\begin{document}
{\sffamily\bfseries\fontsize{23}{27}\selectfont\color{Ink}Safe from Fine-Tune to Phone\par}
\vspace{3pt}
{\sffamily\fontsize{15}{19}\selectfont\color{Accent}Mac experiment report\par}
\vspace{4pt}
{\small\sffamily\color{Muted}Manish | Work completed 4-5 November 2026 | M1-M8\par}
{\footnotesize\sffamily\color{Muted}Prepared 6 November 2026. Evidence and pending checks refer to these recorded experiments.\par}
\vspace{5pt}
{\color{Accent}\hrule height 1.2pt}
\vspace{5pt}
\section{1. What we are trying to find
out}\label{what-we-are-trying-to-find-out}

The project follows Tiny Aya from benign medical fine-tuning to
compressed deployment in a clinic app, in English, Hindi and Hinglish.
The main question is whether safety survives that journey while medical
usefulness is preserved. Later experiments ask whether refusal circuits
can guide a training-free repair and protect that repair without
increasing file size.

\textbf{Current conclusion:} the Mac deployment pipeline works and the
first response files are ready. We have not yet established any safety
improvement, safety loss or successful repair. Those conclusions need
judge labels and the remaining controlled experiments.

\section{2. What I did and what
changed}\label{what-i-did-and-what-changed}

\begin{itemize}
\tightlist
\item
  Set up the Apple M5/24 GB environment, built llama.cpp with Metal, and
  pinned the software versions.
\item
  Downloaded Global and Base, verified their published weight checksums,
  and checked that all 290 tensor names/shapes match.
\item
  Checked evaluation inputs. Each full cell actually has \textbf{856
  harmful + 756 benign = 1,612 prompts}. I retained all records instead
  of trimming to the plan\textquotesingle s rounded 1,600. Eight
  training-package checksum differences are explained by Windows/Linux
  line endings; evaluation checksums pass.
\item
  Merged the medical and control adapters on CPU in FP32 and saved FP16
  checkpoints. Each has 290 finite FP16 tensors. Current cross-machine
  merge parity is still unverified.
\item
  Built original-model F16, Q4\_K\_M and Q3\_K\_M GGUFs, plus the
  medical F16 GGUF. Approximate original file sizes are 6.71 GB, 2.14 GB
  and 1.80 GB.
\item
  Built a resumable AuditKIT-to-llama-server harness. It checks
  model/input/template/code hashes, prompt token IDs and context limits.
  Test generation requires a committed protocol. The Mac protocol and
  input IDs were committed as \textbf{2582b4d before test generation}.
\item
  Verified the official and no-preamble templates. The latter removes
  the system turn while preserving conversation delimiters. Explicit
  token IDs prevent duplicate beginning-of-sequence tokens.
\item
  Built and tested an exporter that redistributes precision within a
  fixed file-size budget, including a matched random control.
\end{itemize}

\section{3. Measured results}\label{measured-results}

{\def\LTcaptype{none} % do not increment counter
\begin{longtable}[]{@{}
  >{\raggedright\arraybackslash}p{(\linewidth - 8\tabcolsep) * \real{0.4000}}
  >{\raggedleft\arraybackslash}p{(\linewidth - 8\tabcolsep) * \real{0.1200}}
  >{\raggedleft\arraybackslash}p{(\linewidth - 8\tabcolsep) * \real{0.1400}}
  >{\raggedleft\arraybackslash}p{(\linewidth - 8\tabcolsep) * \real{0.2100}}
  >{\raggedleft\arraybackslash}p{(\linewidth - 8\tabcolsep) * \real{0.1300}}@{}}
\toprule\noalign{}
\begin{minipage}[b]{\linewidth}\raggedright
Run
\end{minipage} & \begin{minipage}[b]{\linewidth}\raggedleft
Responses
\end{minipage} & \begin{minipage}[b]{\linewidth}\raggedleft
Time
\end{minipage} & \begin{minipage}[b]{\linewidth}\raggedleft
Generated tokens/second
\end{minipage} & \begin{minipage}[b]{\linewidth}\raggedleft
Request failures
\end{minipage} \\
\midrule\noalign{}
\endhead
\bottomrule\noalign{}
\endlastfoot
Original Q4\_K\_M, official template & 1,612 & 59.34 min & 89.87 & 0 \\
Same Q4\_K\_M, no safety preamble & 1,612 & 46.27 min & 111.82 & 0 \\
Medical F16, proposed fidelity set & 50 & 3.23 min & 66.14 & 0 \\
\end{longtable}
}

The hand-off contains \textbf{3,274 responses}. The separate 200-prompt
development benchmark completed at about 103 generated tokens/second; it
is not included in that total.

Both original-model cells use identical prompt IDs, model weights and
controlled generation settings. The full runs cover 554 English, 554
Hindi and 504 Hinglish prompts each. The fidelity set has 18 English, 16
Hindi and 16 Hinglish prompts, with 25 harmful and 25 benign calibration
items.

These numbers measure execution and speed. Zero request failures does
not mean zero unsafe answers. The timing difference is an observation
from these runs, not an isolated or repeated performance experiment:
output token totals also differ.

Generated-token counts include control tokens. In particular, the
medical token rate should not be interpreted as useful visible-answer
throughput; see the termination finding below.

Nine targeted software tests pass. The response
archive\textquotesingle s integrity, file checksums, IDs and run
fingerprints were also verified.
\href{https://github.com/SushovitNanda/Lexsi_Hackathon/tree/manishshaw/results/mac/}{Run
summaries and manifests} contain the underlying evidence.

\section{4. Important findings}\label{important-findings}

\textbf{Stopping configuration matters.} The original
model\textquotesingle s generation config uses end-of-sequence (EOS)
token ID 3, while its tokenizer names token 6 as EOS. The stock
converter selected 6. I normalized every GGUF to the source generation
setting, 3, before the registered test runs. A benign smoke test then
stopped correctly. This change is recorded in the build recipe and
hashes.

\textbf{The medical model has a termination issue to investigate.} All
50 medical responses reached the 256-token cap. In a separate benign
64-token probe, the visible answer was a normal short sentence, followed
by 51 hidden end-of-turn tokens (6), one end-response token, and no EOS
token (3). The training chat template ends assistant turns with 6; the
inherited generation config uses 3. Original/medical GGUF architecture
metadata match.

This supports investigating a training-format/stopping mismatch. It does
not establish the cause across all responses, prove that visible answers
were cut short, or prove a safety loss. Confirm the behavior on the T4.
I retained the frozen outputs and stopping rule; any change for later
comparisons needs a documented protocol amendment.
\href{https://github.com/SushovitNanda/Lexsi_Hackathon/blob/manishshaw/results/mac/medical_termination_diagnostic.json}{Diagnostic
evidence}.

\textbf{Equal-size precision redistribution is feasible.} With synthetic
placeholder scores, eight matrices were upgraded and 39 downgraded.
Protected and random-control files are exactly \textbf{1,795,391,040
bytes}, both \textbf{0.2417\% below stock}, within the required 1\%
tolerance. All 243 unselected tensors remain byte-identical to stock.
The exporter explicitly preserves baseline matrix types because partial
overrides can shift llama.cpp\textquotesingle s default mixture
elsewhere.
\href{https://github.com/SushovitNanda/Lexsi_Hackathon/blob/manishshaw/results/mac/placeholder_protection_report.json}{Exporter
evidence}.

The scores were synthetic. This validates exporter mechanics, not
circuit discovery or safety protection.

\section{5. How to infer meaning from the response
files}\label{how-to-infer-meaning-from-the-response-files}

\begin{enumerate}
\def\labelenumi{\arabic{enumi}.}
\tightlist
\item
  Resolve the current medical/control reference hashes. If they differ,
  measure maximum absolute tensor differences using reference weights; a
  hash difference alone does not measure numerical error.
\item
  Confirm the proposed 50 fidelity prompts, generate the same medical
  FP16 responses on Sushovit\textquotesingle s T4, and judge both
  versions on the same pinned judge host. The planned fidelity criterion
  is at least 48/50 matching judge labels. It has not been evaluated.
\item
  Centrally judge the original cells after judge calibration. Preserve
  raw labels and refusal annotations. Only prompt text and response
  should enter the judge; model/setting metadata stays outside its
  context.
\item
  Join by \textbf{cell ID and prompt ID}, not line order. Calculate
  attack success rate (ASR) as unsafe responses divided by harmful
  prompts. Calculate over-refusal rate (ORR) as refused responses
  divided by benign prompts. Separate language and clinical/general harm
  groups. Controversial counts as safe in primary ASR and unsafe in
  sensitivity analysis.
\item
  Compare no-preamble minus official ASR/ORR using paired prompt IDs. A
  positive ASR difference indicates more unsafe answers in that setting;
  a positive ORR difference indicates more refusals of benign requests.
  Use the registered paired bootstrap and agreed Holm correction before
  claiming a reliable difference.
\end{enumerate}

Neither response length, speed nor successful file conversion
substitutes for these safety measurements.

\clearpage

\section{6. What this means for each research
question}\label{what-this-means-for-each-research-question}

{\def\LTcaptype{none} % do not increment counter
\begin{longtable}[]{@{}
  >{\raggedright\arraybackslash}p{(\linewidth - 4\tabcolsep) * \real{0.2200}}
  >{\raggedright\arraybackslash}p{(\linewidth - 4\tabcolsep) * \real{0.3400}}
  >{\raggedright\arraybackslash}p{(\linewidth - 4\tabcolsep) * \real{0.4400}}@{}}
\toprule\noalign{}
\begin{minipage}[b]{\linewidth}\raggedright
Research question
\end{minipage} & \begin{minipage}[b]{\linewidth}\raggedright
What this work contributes
\end{minipage} & \begin{minipage}[b]{\linewidth}\raggedright
What remains unanswered
\end{minipage} \\
\midrule\noalign{}
\endhead
\bottomrule\noalign{}
\endlastfoot
RQ1: language/clinical safety gaps & Multilingual original Q4 response
data and domain tags. & No judge-labelled gaps yet; Q4 data also do not
establish the FP16 baseline. \\
RQ2: medical fine-tuning drift & Medical/control FP16 merges and a
termination diagnostic. & No matched original/medical/control safety
comparison; the 50 fidelity prompts cannot establish drift. \\
RQ3: refusal circuits & An exporter accepting layer/matrix score tables.
& No validated circuit result from this work; placeholder scores are not
CircuitKIT evidence. \\
RQ4: training-free repair & Conversion, checks and reproducible
evaluation infrastructure. & No repaired checkpoint evaluated here;
repair effectiveness and utility preservation remain open. \\
RQ5: deployment survival & A controlled original-Q4 template comparison,
medical F16 conversion and fidelity responses. & Labels are needed for
template safety effects. Two Q4 cells cannot isolate quantization loss.
Repair retention needs medical/repaired pairs and F16 reference
gains. \\
RQ6: circuit-aware protection & Real protected/random exports at the
same size, with verified tensor integrity. & Safety benefit needs a
validated circuit, selected repaired model and judged
stock/protected/random runs. \\
\end{longtable}
}

For the overall clinic problem, this work makes the deployment
experiment reproducible and exposes a configuration issue that could
otherwise confuse interpretation. The central safety hypotheses remain
open. A smaller file, faster run or working exporter alone does not
demonstrate that a clinic model retains safety or medical usefulness.

\section{7. Limitations and blockers specific to my
experiments}\label{limitations-and-blockers-specific-to-my-experiments}

\textbf{What my part does establish:} the Mac can execute the registered
response runs, the selected files and inputs are traceable, the local
merges are finite FP16, and the precision exporter meets the tested size
budget. The following limits prevent stronger conclusions.

{\def\LTcaptype{none} % do not increment counter
\begin{longtable}[]{@{}
  >{\raggedright\arraybackslash}p{(\linewidth - 4\tabcolsep) * \real{0.2000}}
  >{\raggedright\arraybackslash}p{(\linewidth - 4\tabcolsep) * \real{0.3800}}
  >{\raggedright\arraybackslash}p{(\linewidth - 4\tabcolsep) * \real{0.4200}}@{}}
\toprule\noalign{}
\begin{minipage}[b]{\linewidth}\raggedright
Part of my work
\end{minipage} & \begin{minipage}[b]{\linewidth}\raggedright
Limitation or blocker
\end{minipage} & \begin{minipage}[b]{\linewidth}\raggedright
What it prevents / next action
\end{minipage} \\
\midrule\noalign{}
\endhead
\bottomrule\noalign{}
\endlastfoot
M1-M4: Mac setup and timing & One Apple M5/24 GB machine, one
engine/version and individual timing runs. NF4/CUDA and real phone
hardware were not tested on this Mac. & Do not generalize these speeds
or safety outcomes to T4, Android/iOS hardware or all inference engines.
Keep later comparisons within the registered engine; report
hardware/configuration. \\
M5: adapter merges & All tensors are finite FP16, but local hashes
differ from historical logs. Current D2 reference manifests/tensors have
not been supplied. & Numerical merge parity is unverified. Sushovit must
supply current references; measure actual tensor differences if hashes
differ. A successful merge is not a parity test. \\
M6: safety comparison & We have responses but no centralized judge
labels, calibrated judge results or final shared statistical-family
alignment. The older shared PREREG remains a draft. & ASR, ORR, language
gaps, confidence intervals and significance are not available. Obtain
G1/D7 and align D3 before statistical inference. \\
Evaluation data & Language totals are unequal; per-category samples are
small. Some Hindi/Hinglish sets were drafted or rewritten using models.
My runs do not verify completion of the required native-speaker review.
& Report sample counts and uncertainty. Confirm data review and
per-language judge calibration with Sushovit; distinguish paired
within-language setting comparisons from cross-language claims. \\
M7: engine fidelity & Only 50 proposed calibration prompts were
generated locally, with no agreed replacement set, T4 reference
responses or paired judge labels received. & The gate requires at least
48 matching labels out of 50 and is pending. Passing it would establish
label agreement on this sample, not exact weights, logits or responses
across engines. \\
Medical termination & All 50 medical runs hit the token cap; one benign
probe repeatedly generated hidden token 6 instead of EOS 3. & Review the
same behavior on T4 before further reported medical cells. Any
stopping-rule change must be documented and applied consistently to the
affected comparisons; retain the original outputs. \\
M8: protection & The tested scores are synthetic. No validated D4
circuit or selected repaired checkpoint was used, and no safety
evaluation of these mixed-precision builds was performed. Equal file
size also does not establish equal RAM, latency or energy use. & This is
an exporter feasibility result. Validated circuit/repair hand-offs and
judged stock/protected/random cells are required for a protection
claim. \\
\end{longtable}
}

The \textbf{current external blockers} are D2 merge references,
confirmation of the shared fidelity inputs, T4 responses, G1/D7 judge
labels and D3 statistical alignment. The termination review is a joint
configuration decision. D4 circuit validation and D5 repair selection
are \textbf{future dependencies} for the later RQ4/RQ6 experiments,
rather than unfinished exporter mechanics from today. These are the
limits of my completed local work; none is treated as a passed gate.

\section{8. Handoff and current
status}\label{handoff-and-current-status}

The
\href{https://github.com/SushovitNanda/Lexsi_Hackathon/blob/manishshaw/models/local-results/day2/D6_Manish_2026-11-05.zip}{D6
ZIP} contains responses, fixed fidelity inputs, manifests, the committed
protocol, a T4 helper and judging instructions. It was created as a
private hand-off, then \textbf{pushed to this public repository at the
user\textquotesingle s explicit request}, in commit \textbf{cef0e55}.
Its internal README/manifest preserve the creation-time transfer status.
Repository upload is confirmed; Sushovit\textquotesingle s receipt and
judging are not confirmed. Model weight files remain local.

Still needed: current D2 merge references, agreement on the fidelity
inputs, T4 responses, centralized labels after judge calibration, and
alignment with Sushovit\textquotesingle s final shared pre-registration
before statistical inference. Review the medical termination issue
before further reported medical runs.

The detailed execution record is
\href{https://github.com/SushovitNanda/Lexsi_Hackathon/blob/manishshaw/MAC_NOV5.md}{MAC\_NOV5.md};
the
\href{https://github.com/SushovitNanda/Lexsi_Hackathon/blob/manishshaw/results/mac/day2_handoff.json}{handoff
manifest} links the files to their verified hashes.
\end{document}
